%======================================================================
%  GRAPHRACER -- Presentation Beamer
%  Auteurs : SAAD OMIT & IZADINE
%======================================================================
\documentclass[aspectratio=169,12pt]{beamer}

%----------------------------------------------------------------------
% Theme & couleurs
%----------------------------------------------------------------------
\usetheme{Madrid}
\usecolortheme{default}

\definecolor{darkbg}{HTML}{020617}
\definecolor{mycyan}{HTML}{22d3ee}
\definecolor{myblue}{HTML}{2563eb}
\definecolor{mypurple}{HTML}{9333ea}
\definecolor{mygreen}{HTML}{22c55e}
\definecolor{myred}{HTML}{ef4444}
\definecolor{myyellow}{HTML}{facc15}
\definecolor{myslate}{HTML}{334155}
\definecolor{mylightslate}{HTML}{cbd5e1}

\setbeamercolor{palette primary}{bg=darkbg,fg=white}
\setbeamercolor{palette secondary}{bg=darkbg,fg=mycyan}
\setbeamercolor{palette tertiary}{bg=darkbg,fg=white}
\setbeamercolor{palette quaternary}{bg=darkbg,fg=white}
\setbeamercolor{structure}{fg=mycyan}
\setbeamercolor{background canvas}{bg=darkbg}
\setbeamercolor{normal text}{fg=white}
\setbeamercolor{frametitle}{bg=darkbg,fg=mycyan}
\setbeamercolor{block title}{bg=myblue,fg=white}
\setbeamercolor{block body}{bg=myslate!30!darkbg,fg=white}
\setbeamercolor{block title alerted}{bg=myred,fg=white}
\setbeamercolor{block body alerted}{bg=myred!20!darkbg,fg=white}
\setbeamercolor{block title example}{bg=mygreen,fg=darkbg}
\setbeamercolor{block body example}{bg=mygreen!20!darkbg,fg=white}
\setbeamercolor{title}{fg=mycyan}
\setbeamercolor{subtitle}{fg=mylightslate}
\setbeamercolor{author}{fg=white}
\setbeamercolor{date}{fg=mylightslate}

%----------------------------------------------------------------------
% Packages
%----------------------------------------------------------------------
\usepackage[french]{babel}
\usepackage[utf8]{inputenc}
\usepackage[T1]{fontenc}
\usepackage{listings}
\usepackage{xcolor}
\usepackage{tikz}
\usepackage{pgfplots}
\pgfplotsset{compat=1.18}
\usepackage{hyperref}
\usepackage{array}
\usepackage{booktabs}
\usepackage{dirtree}

%----------------------------------------------------------------------
% Style listings JavaScript (ASCII uniquement dans les blocs de code)
%----------------------------------------------------------------------
\lstdefinelanguage{JavaScript}{
  keywords={const,let,var,function,return,if,else,for,while,new,this,
             export,import,from,of,class,extends,async,await,true,false,null},
  keywordstyle=\color{mycyan}\bfseries,
  ndkeywords={useState,useEffect,useRef,useCallback,useContext,createContext},
  ndkeywordstyle=\color{myyellow},
  identifierstyle=\color{white},
  sensitive=false,
  comment=[l]{//},
  morecomment=[s]{/*}{*/},
  commentstyle=\color{mylightslate}\itshape,
  stringstyle=\color{mygreen},
  morestring=[b]',
  morestring=[b]",
  morestring=[b]`
}

\lstset{
  language=JavaScript,
  basicstyle=\ttfamily\footnotesize\color{white},
  backgroundcolor=\color{myslate!20!darkbg},
  frame=single,
  framerule=0.5pt,
  rulecolor=\color{myslate},
  breaklines=true,
  showstringspaces=false,
  tabsize=2,
  xleftmargin=0.5em,
  xrightmargin=0.5em,
  lineskip=-0.5pt,
  extendedchars=false,
  inputencoding=utf8
}

%----------------------------------------------------------------------
% Macros
%----------------------------------------------------------------------
\newcommand{\tech}[1]{\textcolor{mycyan}{\textbf{\texttt{#1}}}}
\newcommand{\highlight}[1]{\textcolor{myyellow}{\textbf{#1}}}
\newcommand{\success}[1]{\textcolor{mygreen}{\textbf{#1}}}
\newcommand{\danger}[1]{\textcolor{myred}{\textbf{#1}}}

%======================================================================
\title{\Huge\textbf{GRAPH}\textcolor{white}{RACER}}
\subtitle{Engineer Edition --- Pr\'esentation Technique \& Business Plan}
\author{\textbf{SAAD OMIT} \and \textbf{IZADINE}}
\date{\today}
%======================================================================

\begin{document}

%---------- SLIDE TITRE ------------------------------------------------
\begin{frame}[plain]
  \vfill
  \centering
  {\Huge \textcolor{mycyan}{\textbf{GRAPH}}\textbf{\textcolor{white}{RACER}}}\\[0.4cm]
  {\large \textcolor{mylightslate}{Engineer Edition}}\\[1cm]
  \begin{tikzpicture}
    \draw[mycyan, thick, ->] (0,0) -- (8,0) node[right,white] {$x$};
    \draw[mycyan, thick, ->] (0,-0.5) -- (0,2.5) node[above,white] {$f(x)$};
    \draw[myyellow, very thick, domain=0:7.5, samples=80]
      plot (\x, {1.2*sin(\x r) + 1.2});
    \fill[myyellow] (4.8, 1.85) rectangle (5.6, 2.25);
    \fill[darkbg!80!white] (5.0, 2.25) rectangle (5.4, 2.5);
    \fill[darkbg] (5.0, 1.7) circle (0.15);
    \fill[darkbg] (5.6, 1.7) circle (0.15);
    \node[myyellow] at (2.5, 2.6) {\large $\bigstar$};
  \end{tikzpicture}\\[0.8cm]
  {\large \textbf{SAAD OMIT} \quad \& \quad \textbf{IZADINE}}\\[0.3cm]
  {\small \textcolor{mylightslate}{\today}}
  \vfill
\end{frame}

%---------- PLAN -------------------------------------------------------
\begin{frame}{Plan de la Pr\'esentation}
  \tableofcontents
\end{frame}

%======================================================================
\section{Pr\'esentation du Projet}
%======================================================================

\begin{frame}{Qu'est-ce que GraphRacer ?}
  \begin{columns}
    \column{0.55\textwidth}
    \begin{block}{Concept}
      GraphRacer est un \highlight{jeu \'educatif de course} o\`u le joueur pilote
      une voiture en \'ecrivant des \textbf{\'equations math\'ematiques}.
      La courbe trac\'ee devient la route.
    \end{block}
    \vspace{0.4cm}
    \begin{itemize}
      \item \textbf{Public cible :} Lyc\'eens / \'Etudiants (ITS)
      \item \textbf{Mode solo :} 8 niveaux progressifs
      \item \textbf{Mode multi :} Course via QR Code + t\'el\'ephone
      \item \textbf{Objectif :} Apprendre les fonctions maths en jouant
    \end{itemize}
    \column{0.42\textwidth}
    \begin{exampleblock}{Fonctions utilisables}
      \begin{itemize}
        \item Lin\'eaire : \tech{0.5 * x}
        \item Parabole : \tech{-0.1*(x-20)\^{}2+10}
        \item Sinus : \tech{5*sin(x/5)}
        \item Valeur absolue : \tech{-abs(x-20)+15}
        \item Racine carr\'ee : \tech{sqrt(x)}
        \item Combinaisons !
      \end{itemize}
    \end{exampleblock}
  \end{columns}
\end{frame}

%----------------------------------------------------------------------
\begin{frame}{Stack Technologique}
  \begin{columns}
    \column{0.48\textwidth}
    \begin{block}{Front-end}
      \begin{itemize}
        \item \tech{React 19} --- SPA avec routing
        \item \tech{React Router DOM} --- Navigation
        \item \tech{Canvas API} --- Moteur graphique 2D
        \item \tech{Web Audio API} --- Sons synth\'etis\'es
        \item \tech{Tailwind CSS} --- Styles
        \item \tech{lucide-react} --- Ic\^ones
        \item \tech{qrcode.react} --- G\'en\'eration QR
      \end{itemize}
    \end{block}
    \column{0.48\textwidth}
    \begin{block}{Back-end}
      \begin{itemize}
        \item \tech{Node.js + Express} --- Serveur HTTP
        \item \tech{Socket.io} --- Temps r\'eel (WebSocket)
      \end{itemize}
    \end{block}
    \vspace{0.3cm}
    \begin{alertblock}{Architecture}
      Application \textbf{full-stack JavaScript} : m\^eme langage partout.
      Le serveur ne g\`ere que la \textbf{signalisation} --- la physique tourne
      enti\`erement \textbf{c\^ot\'e client}.
    \end{alertblock}
  \end{columns}
\end{frame}

%======================================================================
\section{Architecture du Code}
%======================================================================

\begin{frame}{Structure des Fichiers}
  \begin{columns}
    \column{0.48\textwidth}
    \dirtree{%
      .1 graphracer/.
      .2 \textcolor{mycyan}{server.js}.
      .2 package.json.
      .2 src/.
      .3 \textcolor{mycyan}{App.js}.
      .3 \textcolor{myyellow}{context/}.
      .4 GameContext.jsx.
      .3 \textcolor{myyellow}{engine/}.
      .4 MathParser.js.
      .3 \textcolor{myyellow}{data/}.
      .4 levels.js.
      .3 \textcolor{myyellow}{pages/}.
      .4 Home.jsx.
      .4 LevelSelect.jsx.
      .4 \textcolor{mycyan}{Game.jsx}.
      .4 Garage.jsx.
      .4 Settings.jsx.
      .4 \textcolor{mycyan}{MultiplayerHost.jsx}.
      .4 \textcolor{mycyan}{RemoteController.jsx}.
    }
    \column{0.48\textwidth}
    \begin{tikzpicture}[scale=0.75, every node/.style={font=\small}]
      \node[draw=mycyan,fill=myslate!30!darkbg,text=white,rounded corners,minimum width=3cm,minimum height=0.6cm] (app) at (3,5) {App.js (Router)};
      \node[draw=myyellow,fill=myslate!30!darkbg,text=white,rounded corners,minimum width=2.5cm] (ctx) at (1,3.5) {GameContext};
      \node[draw=mycyan,fill=myslate!30!darkbg,text=white,rounded corners,minimum width=2.5cm] (game) at (4.5,2) {Game.jsx};
      \node[draw=mycyan,fill=myslate!30!darkbg,text=white,rounded corners,minimum width=2.5cm] (mhost) at (1,2) {MultiHost};
      \node[draw=mygreen,fill=myslate!30!darkbg,text=white,rounded corners,minimum width=2.5cm] (math) at (4.5,0.5) {MathParser};
      \node[draw=mypurple,fill=myslate!30!darkbg,text=white,rounded corners,minimum width=2.5cm] (srv) at (1,0.5) {server.js};
      \node[draw=mypurple,fill=myslate!30!darkbg,text=white,rounded corners,minimum width=2.5cm] (rc) at (4.5,3.5) {RemoteCtrl};
      \draw[->,mycyan] (app) -- (ctx);
      \draw[->,mycyan] (app) -- (game);
      \draw[->,mycyan] (app) -- (mhost);
      \draw[->,mycyan] (app) -- (rc);
      \draw[->,mygreen] (game) -- (math);
      \draw[->,mypurple,dashed] (mhost) -- (srv);
      \draw[->,mypurple,dashed] (rc) -- (srv);
    \end{tikzpicture}
  \end{columns}
\end{frame}

%----------------------------------------------------------------------
\begin{frame}{Routing et Navigation (App.js)}
  \begin{block}{7 routes React Router DOM}
    \vspace{0.2cm}
    \begin{tabular}{lll}
      \toprule
      \textcolor{mycyan}{Route} & \textcolor{myyellow}{Composant} & \textcolor{mylightslate}{Description}\\
      \midrule
      \texttt{/} & \tech{Home} & Menu principal (1J / Multi / Garage)\\
      \texttt{/select} & \tech{LevelSelect} & Grille des 8 niveaux\\
      \texttt{/game} & \tech{Game} & Moteur de jeu solo\\
      \texttt{/garage} & \tech{Garage} & Personnalisation couleurs\\
      \texttt{/settings} & \tech{Settings} & Options (grille\ldots)\\
      \texttt{/multi-host} & \tech{MultiplayerHost} & PC h\^ote course\\
      \texttt{/join} & \tech{RemoteController} & T\'el\'ephone joueur\\
      \bottomrule
    \end{tabular}
  \end{block}
  \vspace{0.3cm}
  L'\'etat global (niveau, couleurs, progression) transite par \tech{GameContext}
  via le pattern \textbf{React Context + localStorage}.
\end{frame}

%----------------------------------------------------------------------
\begin{frame}{\'Etat Global --- GameContext.jsx}
  \begin{columns}
    \column{0.5\textwidth}
    \begin{block}{\'Etats persist\'es (localStorage)}
      \begin{itemize}
        \item \tech{currentLevel} --- niveau actif
        \item \tech{unlockedLevels[]} --- niveaux d\'ebloqués
        \item \tech{carColor} --- couleur voiture
        \item \tech{lineColor} --- couleur trac\'e
        \item \tech{bgColor} --- fond de jeu
        \item \tech{showGrid} --- affichage grille
      \end{itemize}
    \end{block}
    \column{0.48\textwidth}
\begin{lstlisting}
// Auto-save to localStorage
useEffect(() => {
  localStorage.setItem(
    'carColor', carColor);
  localStorage.setItem(
    'unlockedLevels',
    JSON.stringify(unlockedLevels));
}, [carColor, unlockedLevels]);

// Progressive unlock
const unlockNextLevel = (id) => {
  if(!unlockedLevels.includes(id+1))
    setUnlockedLevels(
      prev => [...prev, id+1]);
};
\end{lstlisting}
  \end{columns}
\end{frame}

%======================================================================
\section{Fonctionnalit\'es Impl\'ement\'ees}
%======================================================================

\begin{frame}{Fonctionnalit\'e 1 --- Moteur Math\'ematique (MathParser.js)}
  \begin{block}{Principe : \'Equation texte $\rightarrow$ Points $\rightarrow$ Courbe}
  \end{block}
\begin{lstlisting}
export const generateCurve = (equationStr, distance) => {
  // 1. Nettoyage : remplace les mots par du JS pur
  let safeEq = equationStr.toLowerCase()
    .replace(/sin/g, 'Math.sin')
    .replace(/abs/g, 'Math.abs')
    .replace(/\^/g,  '**');   // x^2 -> x**2

  // 2. Compilation dynamique
  const f = new Function('x', `return ${safeEq}`);

  // 3. Discretisation (pas de 0.1 metre)
  for (let x = 0; x <= distance; x += 0.1) {
    const y = -f(x);  // inversion Y (repere ecran)
    if (!isFinite(y) || isNaN(y)) continue;
    points.push({ x: x*SCALE, y: y*SCALE, realX: x, realY: -y });
  }
};
\end{lstlisting}
  \vspace{0.2cm}
  \textcolor{myyellow}{\textbf{Note :}} La constante \tech{SCALE = 40} convertit les m\`etres en pixels.
  L'inversion de $Y$ est indispensable car l'axe Y descend en informatique.
\end{frame}

%----------------------------------------------------------------------
\begin{frame}{Fonctionnalit\'e 2 --- Tron\c{c}ons de Fonctions}
  \begin{columns}
    \column{0.5\textwidth}
    \begin{block}{Idée}
      Le joueur peut d\'efinir \textbf{plusieurs segments} de fonctions
      sur des intervalles distincts.\\[0.4cm]
      Exemple niveau 6 « Les Tron\c{c}ons » :\\
      \begin{enumerate}
        \item $f(x) = 0$ sur $[0 ; 30]$
        \item $f(x) = x - 30$ sur $[30 ; 70]$
      \end{enumerate}
    \end{block}
    \column{0.48\textwidth}
\begin{lstlisting}
// Calcul Y global a targetX
const getWorldY = (targetX) => {
 // Trouve le segment actif
 const seg = [...segments]
   .reverse()
   .find(s =>
     targetX >= s.start &&
     targetX <  s.end);
 if (!seg) return null;

 // Continuite entre segments
 let startH = 0;
 if (seg.start > 0)
  startH = getWorldY(seg.start-0.01);

 const localX = targetX - seg.start;
 return startH + eval(seg.eq, localX);
};
\end{lstlisting}
  \end{columns}
\end{frame}

%----------------------------------------------------------------------
\begin{frame}{Fonctionnalit\'e 3 --- Moteur de Jeu (Game.jsx)}
  \begin{block}{Double Canvas pour les performances}
    \tech{bgCanvasRef} : sol statique, redessin\'e \`a chaque changement d'\'equation\\
    \tech{carCanvasRef} : voiture + particules, redessin\'e \`a 60 FPS
  \end{block}
\begin{lstlisting}
// Boucle RAF (requestAnimationFrame)
const gameLoop = useCallback(() => {
  // 1. Physique particules (fumee, explosion)
  gameState.current.particles.forEach(p => {
    p.x += p.vx;  p.y += p.vy;
    p.life -= p.decay;
  });

  // 2. Deplacement + consommation carburant
  if (autoDrive) {
    gameState.current.x    += 0.3;
    gameState.current.fuel -= 0.15;
  }

  // 3. Detection collision (trou / obstacle)
  if (currentY < -5) setGameStatus('crashed');

  // 4. Dessin + boucle suivante
  drawCarFrame();
  requestRef.current = requestAnimationFrame(gameLoop);
}, [gameStatus, currentLevel, drawCarFrame]);
\end{lstlisting}
\end{frame}

%----------------------------------------------------------------------
\begin{frame}{Fonctionnalit\'e 3 (suite) --- Physique \& Collisions}
  \begin{columns}
    \column{0.48\textwidth}
    \begin{block}{Gestion des obstacles}
\begin{lstlisting}
// Trou dans le sol
if (currentY < -5)
  crashed = true;

// Mur rectangulaire
obstacles.forEach(obs => {
  if (nextX >= obs.x &&
      nextX <= obs.x+obs.width &&
      currentY < obs.y+obs.height)
    crashed = true;
});

// Detection etoile
const dist = Math.sqrt(
  (nextX  - star.x)**2 +
  (currentY - star.y)**2);
if (dist < 2.5) collectStar(idx);
\end{lstlisting}
    \end{block}
    \column{0.48\textwidth}
    \begin{block}{Effets visuels}
      \begin{itemize}
        \item \success{Fum\'ee} grise derri\`ere la voiture (particules)
        \item \danger{Explosion} rouge/orange au crash (30 particules)
        \item \textcolor{myyellow}{\'Etoiles} pulsantes (animation \texttt{sin(time)})
        \item \highlight{Inclinaison} selon la d\'eriv\'ee (\texttt{atan2})
        \item \textcolor{mycyan}{Flamme} r\'eacteur (cercle orange anim\'e)
      \end{itemize}
    \end{block}
    \begin{exampleblock}{Angle de la voiture}
      \small $\theta = \arctan\!\left(\dfrac{\Delta Y}{\Delta X}\right)$
      calcul\'e entre deux points cons\'ecutifs
    \end{exampleblock}
  \end{columns}
\end{frame}

%----------------------------------------------------------------------
\begin{frame}{Fonctionnalit\'e 4 --- Audio Synth\'etis\'e (Web Audio API)}
  \begin{columns}
    \column{0.5\textwidth}
\begin{lstlisting}
// Son victoire (accord Do Mi Sol)
const playWinSound = () => {
  const ctx = new AudioContext();
  const now = ctx.currentTime;
  [523.25, 659.25, 783.99]
    .forEach((freq, i) => {
      const osc = ctx.createOscillator();
      const gain = ctx.createGain();
      osc.type = 'sine';
      osc.frequency.value = freq;
      gain.gain.setValueAtTime(0.1, now);
      gain.gain.exponentialRampToValueAtTime(
        0.001, now + 0.5 + i*0.1);
      osc.start(now + i*0.1);
    });
};
\end{lstlisting}
    \column{0.48\textwidth}
    \begin{block}{3 sons diff\'erents}
      \begin{itemize}
        \item \success{Victoire} --- accord Do-Mi-Sol\\
          (oscillateur \texttt{sine})
        \item \textcolor{myyellow}{\'Etoile} --- son aigu montant\\
          (oscillateur \texttt{triangle})
        \item \danger{Crash} --- son grave descendant\\
          (oscillateur \texttt{sawtooth})
      \end{itemize}
    \end{block}
    \vspace{0.3cm}
    \textbf{Z\'ero fichier audio !} Tout est g\'en\'er\'e \`a la vol\'ee par
    l'API Web Audio avec des oscillateurs.
  \end{columns}
\end{frame}

%----------------------------------------------------------------------
\begin{frame}{Fonctionnalit\'e 5 --- Syst\`eme de Niveaux (levels.js)}
  \begin{block}{8 niveaux progressifs --- du Tutorial au Cauchemar}
  \end{block}
  \begin{tabular}{clll}
    \toprule
    \textcolor{mycyan}{\#} & \textcolor{myyellow}{Titre} & \textcolor{mylightslate}{Difficult\'e} & \textcolor{white}{Concept math\'ematique}\\
    \midrule
    1 & D\'ecollage & Tuto & Fonction lin\'eaire $ax$\\
    2 & Le Grand Saut & Facile & Saut par-dessus un trou\\
    3 & La Parabole & Moyen & Parabole $-a(x-b)^2+c$\\
    4 & Sinus Valley & Difficile & Fonctions trigonom\'etriques\\
    5 & La Pyramide & Moyen & Valeur absolue $|x|$\\
    6 & Les Tron\c{c}ons & Difficile & Fonctions par morceaux\\
    7 & Slalom G\'eant & Expert & Combinaison $|\sin(x)|$\\
    8 & L'Examen Final & Cauchemar & Tout combiner\\
    \bottomrule
  \end{tabular}
  \vspace{0.2cm}\\
  Chaque niveau d\'efinit : \tech{holes[]} (trous), \tech{obstacles[]} (murs),
  \tech{stars[]} (\'etoiles), \tech{fuelObjective} et un \tech{memo} p\'edagogique.
\end{frame}

%----------------------------------------------------------------------
\begin{frame}{Fonctionnalit\'e 6 --- Mode Multijoueur (Socket.io)}
  \begin{block}{Architecture R\'eseau}
    \begin{center}
    \begin{tikzpicture}[scale=0.72, every node/.style={font=\small}]
      \node[draw=mycyan,fill=myslate!30!darkbg,text=white,rounded corners,
            minimum width=3cm,minimum height=0.8cm] (pc) at (4,3)
            {PC H\^ote (MultiplayerHost)};
      \node[draw=mypurple,fill=myslate!30!darkbg,text=white,rounded corners,
            minimum width=3cm,minimum height=0.8cm] (srv) at (4,1.5)
            {server.js :4000};
      \node[draw=myyellow,fill=myslate!30!darkbg,text=white,rounded corners,
            minimum width=2.2cm] (p1) at (0.5,0) {Joueur 1};
      \node[draw=mygreen,fill=myslate!30!darkbg,text=white,rounded corners,
            minimum width=2.2cm] (p2) at (4,0) {Joueur 2};
      \node[draw=myred,fill=myslate!30!darkbg,text=white,rounded corners,
            minimum width=2.2cm] (p3) at (7.5,0) {Joueur 3};
      \draw[<->,mycyan,thick] (pc) -- node[right,white,font=\tiny]{Socket.io WS} (srv);
      \draw[<->,myyellow,dashed] (p1) -- (srv);
      \draw[<->,mygreen,dashed] (p2) -- (srv);
      \draw[<->,myred,dashed] (p3) -- (srv);
      \node[draw=white,rounded corners,fill=darkbg,text=white,font=\tiny,
            minimum width=3cm] at (4,4.2) {QR Code : http://IP:3000/join};
    \end{tikzpicture}
    \end{center}
  \end{block}
\end{frame}

%----------------------------------------------------------------------
\begin{frame}{Fonctionnalit\'e 6 (suite) --- Flux Multijoueur}
  \begin{columns}
    \column{0.5\textwidth}
    \begin{block}{Serveur (server.js) --- 4 \'ev\'enements}
\begin{lstlisting}
// 1. Joueur rejoint
socket.on('join_game', (data) => {
  players.push({id: socket.id, ...data});
  io.emit('players_update', players);
});

// 2. Joueur envoie ses equations
socket.on('submit_functions', (d) => {
  socket.broadcast.emit(
    'receive_functions', {
      playerId: socket.id,
      segments: d.segments
  });
});

// 3. Hote change le niveau
socket.on('host_update', (data) => {
  socket.broadcast.emit('host_update', data);
});

// 4. Deconnexion
socket.on('disconnect', () => {
  players = players.filter(
    p => p.id !== socket.id);
  io.emit('players_update', players);
});
\end{lstlisting}
    \end{block}
    \column{0.46\textwidth}
    \begin{exampleblock}{Phases de Course}
      \begin{enumerate}
        \item \textbf{lobby} --- QR Code affich\'e
        \item \textbf{waiting} --- chacun tape ses \'equations
        \item \textbf{racing} --- simulation sur le PC
        \item \textbf{finished} --- R\'esultats + historique
      \end{enumerate}
    \end{exampleblock}
    \vspace{0.3cm}
    \begin{alertblock}{Cl\'e de conception}
      Le serveur ne simule \textbf{pas} la physique.
      Il ne fait que \textbf{relayer} les \'ev\'enements.
      Le PC h\^ote recalcule tous les trajets localement.
    \end{alertblock}
  \end{columns}
\end{frame}

%----------------------------------------------------------------------
\begin{frame}{Fonctionnalit\'e 7 --- Garage \& Personnalisation}
  \begin{columns}
    \column{0.5\textwidth}
    \begin{block}{Options disponibles}
      \begin{itemize}
        \item Couleur de la voiture (8 choix)
        \item Couleur du trac\'e math\'ematique
        \item Couleur de l'arri\`ere-plan
        \item Affichage / masquage de la grille
        \item Remise \`a z\'ero de la progression
      \end{itemize}
    \end{block}
    \vspace{0.3cm}
    \begin{block}{Sauvegarde automatique}
      Utilisation du \tech{localStorage} pour persister les
      pr\'ef\'erences sans base de donn\'ees.
    \end{block}
    \column{0.48\textwidth}
\begin{lstlisting}
// Garage.jsx
const COLORS = [
  '#facc15', // Jaune
  '#ef4444', // Rouge
  '#3b82f6', // Bleu
  '#10b981', // Vert
  '#a855f7', // Violet
  '#ec4899', // Rose
  '#ffffff', // Blanc
  '#000000', // Noir
];

// Bouton couleur + selection
{COLORS.map(c => (
  <button
    onClick={() => setCarColor(c)}
    style={{
      background: c,
      border: carColor === c
        ? '3px solid white'
        : 'none'
    }} />
))}
\end{lstlisting}
  \end{columns}
\end{frame}

%======================================================================
\section{Diagrammes et Architecture}
%======================================================================

\begin{frame}{Diagramme de Flux --- Mode Solo}
  \begin{center}
  \begin{tikzpicture}[scale=0.8,
    box/.style={draw=mycyan,fill=myslate!30!darkbg,text=white,rounded corners,
                minimum width=2.5cm,minimum height=0.7cm,font=\small},
    decision/.style={draw=myyellow,fill=myslate!30!darkbg,text=white,diamond,
                     aspect=2.5,font=\scriptsize},
    arr/.style={->,mycyan,thick}]

    \node[box] (home) at (0,5) {Home};
    \node[box] (sel) at (3,5) {LevelSelect};
    \node[box] (game) at (6,5) {Game.jsx};
    \node[box] (eq) at (6,3.5) {Saisie f(x)};
    \node[box] (calc) at (6,2) {calculateTerrain};
    \node[box] (loop) at (6,0.5) {gameLoop (RAF)};
    \node[decision] (crash) at (3,0.5) {crash?};
    \node[decision] (win) at (9,0.5) {arriv\'ee?};
    \node[box,draw=mygreen] (won) at (9,2) {\success{VICTOIRE}};
    \node[box,draw=myred] (lost) at (3,2) {\danger{CRASH}};

    \draw[arr] (home) -- node[above,white,font=\tiny]{JOUER} (sel);
    \draw[arr] (sel) -- node[above,white,font=\tiny]{niveau} (game);
    \draw[arr] (game) -- (eq);
    \draw[arr] (eq) -- node[right,white,font=\tiny]{GO} (calc);
    \draw[arr] (calc) -- (loop);
    \draw[arr] (loop) -- (crash);
    \draw[arr] (loop) -- (win);
    \draw[arr] (crash) -- node[left,myred,font=\tiny]{oui} (lost);
    \draw[arr] (win) -- node[right,mygreen,font=\tiny]{oui} (won);
    \draw[->,mycyan,dashed] (lost) -- node[above,white,font=\tiny]{retry} (eq);
  \end{tikzpicture}
  \end{center}
\end{frame}

%----------------------------------------------------------------------
\begin{frame}{Diagramme de Flux --- Mode Multijoueur}
  \begin{center}
  \begin{tikzpicture}[scale=0.78,
    box/.style={draw=mypurple,fill=myslate!30!darkbg,text=white,rounded corners,
                minimum width=2.8cm,minimum height=0.7cm,font=\small},
    arr/.style={->,mypurple,thick}]

    \node[box] (lobby) at (4,5) {Lobby + QR Code};
    \node[box,draw=myyellow] (join) at (0.5,3.5) {Joueur rejoint};
    \node[box] (wait) at (4,3.5) {En attente fonctions};
    \node[box,draw=myyellow] (send) at (0.5,2) {Envoie f(x)};
    \node[box] (race) at (4,2) {COURSE (racing)};
    \node[box,draw=mygreen] (fin) at (4,0.5) {\success{R\'esultats + Historique}};
    \node[box,draw=mycyan] (rev) at (7.5,2) {\tech{REVANCHE}};

    \draw[arr] (lobby) -- node[right,white,font=\tiny]{PR\'EPARER} (wait);
    \draw[->,myyellow,thick] (join) -- node[above,white,font=\tiny]{join\_game} (wait);
    \draw[->,myyellow,thick] (send) -- node[above,white,font=\tiny]{submit} (race);
    \draw[arr] (wait) -- node[right,white,font=\tiny]{tous pr\^ets} (race);
    \draw[arr] (race) -- (fin);
    \draw[->,mycyan,thick] (fin) -- (rev);
    \draw[->,mycyan,thick,dashed] (rev) -- (wait);
    \draw[->,mypurple,dashed] (fin) to[bend right=40] (lobby);
  \end{tikzpicture}
  \end{center}
\end{frame}

%======================================================================
\section{D\'emonstration de Code Cl\'e}
%======================================================================

\begin{frame}{Code Cl\'e --- Calcul de la Trajectoire Multi-Segments}
\begin{lstlisting}
// Dans Game.jsx -- calculateTerrain()
const getWorldY = (targetX) => {
    const memoKey = Math.round(targetX * 10);
    if (memo[memoKey] !== undefined) return memo[memoKey]; // Cache

    // Trouver le bon segment
    const activeSeg = [...segments].reverse()
        .find(s => targetX >= s.start && targetX < s.end);
    if (!activeSeg) return null; // Pas de fonction -> pas de route

    // Continuite verticale entre segments
    let startHeight = 0;
    if (activeSeg.start > 0) {
        const prevY = getWorldY(activeSeg.start - 0.01);
        startHeight = prevY === null ? 0 : prevY;
    }

    // x local : repart de 0 au debut de chaque segment
    const localX = targetX - activeSeg.start;
    const finalY = startHeight + evaluate(activeSeg.eq, localX);
    memo[memoKey] = finalY; // Memoisation pour la performance
    return finalY;
};
\end{lstlisting}
  \vspace{0.2cm}
  \textcolor{myyellow}{\textbf{Astuce :}} M\'emo\"isation via dictionnaire \texttt{memo} pour \'eviter
  le recalcul $\rightarrow$ indispensable pour les grandes distances.
\end{frame}

%----------------------------------------------------------------------
\begin{frame}{Code Cl\'e --- Dessin de la Voiture (Canvas 2D)}
\begin{lstlisting}
// Dans Game.jsx -- drawCarFrame()
// 1. Angle d'inclinaison (derivee numerique)
const dy = -(nextPoint.realY - point.realY) * scale;
const dx =  (nextPoint.realX - point.realX) * scale;
const angle = Math.atan2(dy, dx);

// 2. Dessin de la voiture
ctx.save();
ctx.translate(px, py);
ctx.rotate(angle);        // Inclinaison selon la pente

// Chassis (arrondi)
ctx.fillStyle = carColor;
ctx.beginPath();
ctx.roundRect(-20, -12, 40, 12, 4);
ctx.fill();

// Cockpit (vitre bleue)
ctx.fillStyle = '#0ea5e9';
ctx.roundRect(-5, -22, 20, 10, [8, 8, 0, 0]);
ctx.fill();

// Roues
ctx.arc(-12, 4, 8, 0, Math.PI*2); // Roue AR
ctx.arc( 12, 4, 8, 0, Math.PI*2); // Roue AV
ctx.restore();
\end{lstlisting}
\end{frame}

%======================================================================
\section{Business Plan}
%======================================================================

\begin{frame}{Business Plan --- Vision \& March\'e}
  \begin{columns}
    \column{0.5\textwidth}
    \begin{block}{Proposition de Valeur}
      GraphRacer rend l'apprentissage des fonctions math\'ematiques
      \textbf{engageant} via la \textbf{gamification} et la \textbf{comp\'etition}.
      Le joueur \textit{ressent} la fonction avant de la comprendre formellement.
    \end{block}
    \vspace{0.3cm}
    \begin{block}{March\'es Cibles}
      \begin{itemize}
        \item \textbf{B2C} : Lyc\'eens, \'etudiants pr\'epa / IUT
        \item \textbf{B2B} : \'Etablissements scolaires, ENT
        \item \textbf{B2G} : Appels d'offres \'Education Nationale
      \end{itemize}
    \end{block}
    \column{0.46\textwidth}
    \begin{exampleblock}{Chiffres Cl\'es EdTech (France)}
      \begin{itemize}
        \item \textbf{3,2 Md\$} March\'e EdTech Europe 2024
        \item \textbf{+15\%/an} croissance post-Covid
        \item \textbf{12 M} \'el\`eves lyc\'ee/coll\`ege
        \item \textbf{68\%} enseignants favorables aux outils gamifi\'es
      \end{itemize}
    \end{exampleblock}
  \end{columns}
\end{frame}

%----------------------------------------------------------------------
\begin{frame}{Business Plan --- Mod\`ele \'Economique}
  \begin{block}{3 Mod\`eles de Revenus}
  \end{block}
  \begin{tabular}{p{2.5cm}p{3.5cm}p{3.5cm}}
    \toprule
    \textcolor{mycyan}{\textbf{Mod\`ele}} & \textcolor{myyellow}{\textbf{Offre}} & \textcolor{mylightslate}{\textbf{Prix estim\'e}}\\
    \midrule
    \success{Freemium} &
    Niveaux 1--4 gratuits, 5--8 payants &
    \textbf{4,99 \euro/mois}\\[4pt]
    \success{Licence \'Ecole} &
    Acc\`es illimit\'e + tableau de bord prof &
    \textbf{500--2000 \euro/an}\\[4pt]
    \success{Marque Blanche} &
    Int\'egration ENT (Pronote, Moodle) &
    \textbf{Devis 10--50 k\euro}\\
    \bottomrule
  \end{tabular}
  \vspace{0.4cm}
  \begin{alertblock}{Avantage concurrentiel}
    Le mode \textbf{multijoueur local sans inscription} (QR Code) est unique :
    d\'eploiement imm\'ediat en classe sans infrastructure IT complexe.
  \end{alertblock}
\end{frame}

%----------------------------------------------------------------------
\begin{frame}{Business Plan --- Feuille de Route}
  \begin{tabular}{p{2cm}p{8cm}}
    \toprule
    \textcolor{mycyan}{\textbf{Phase}} & \textcolor{white}{\textbf{Objectifs}}\\
    \midrule
    \textbf{V1 (actuel)} &
    8 niveaux, mode solo, multijoueur LAN, garage\\[4pt]
    \textbf{V2 (6 mois)} &
    \'Editeur de niveaux, tableau de bord enseignant,
    authentification, scores en ligne\\[4pt]
    \textbf{V3 (12 mois)} &
    Battle Royale (16 joueurs), statistiques par fonction,
    int\'egration LMS, PWA mobile\\[4pt]
    \textbf{V4 (18 mois)} &
    IA adaptative (niveau selon erreurs), calcul int\'egral / diff\'erentiel\\
    \bottomrule
  \end{tabular}
\end{frame}

%----------------------------------------------------------------------
\begin{frame}{Business Plan --- Finances Pr\'evisionnelles}
  \begin{columns}
    \column{0.5\textwidth}
    \begin{block}{Investissement Initial}
      \begin{itemize}
        \item D\'eveloppement V2 : \textbf{15 000 \euro}
        \item Design UX/UI : \textbf{5 000 \euro}
        \item Infrastructure Cloud : \textbf{1 200 \euro/an}
        \item Marketing : \textbf{3 000 \euro}
        \item \textbf{Total Y1 : $\approx$ 24 200 \euro}
      \end{itemize}
    \end{block}
    \column{0.46\textwidth}
    \begin{exampleblock}{Projections Revenus}
      \begin{tabular}{lrr}
        \toprule
        & \textbf{An 1} & \textbf{An 3}\\
        \midrule
        Freemium & 5 k\euro & 80 k\euro\\
        Licences & 20 k\euro & 200 k\euro\\
        Marque B. & 10 k\euro & 100 k\euro\\
        \midrule
        \textbf{Total} & \textbf{35 k} & \textbf{380 k}\\
        \bottomrule
      \end{tabular}
    \end{exampleblock}
    \begin{alertblock}{Seuil de rentabilit\'e}
      \textbf{Estim\'e \`a 14 mois} avec 40 licences \'ecoles + 500 abonn\'es
    \end{alertblock}
  \end{columns}
\end{frame}

%----------------------------------------------------------------------
\begin{frame}{Business Plan --- Analyse SWOT}
  \begin{columns}
    \column{0.5\textwidth}
    \begin{exampleblock}{\success{Forces}}
      \begin{itemize}
        \item Concept original (\'equation = route)
        \item Multi sans infrastructure (QR)
        \item Stack JS moderne, maintenable
        \item Contenu p\'edagogique int\'egr\'e (m\'emos)
        \item 0 d\'ependance lourde (pas de moteur 3D)
      \end{itemize}
    \end{exampleblock}
    \begin{alertblock}{\danger{Faiblesses}}
      \begin{itemize}
        \item Multi limit\'e au LAN actuellement
        \item Pas d'authentification ni BDD
        \item Physique simplifi\'ee (pas de gravit\'e)
      \end{itemize}
    \end{alertblock}
    \column{0.46\textwidth}
    \begin{block}{\highlight{Opportunit\'es}}
      \begin{itemize}
        \item Plan num\'erique \'Education Nationale
        \item Partenariat \'editeurs manuels scolaires
        \item Concours d'innovation p\'edagogique
        \item Export international
      \end{itemize}
    \end{block}
    \begin{block}{\textcolor{mylightslate}{Menaces}}
      \begin{itemize}
        \item Concurrence Desmos, GeoGebra (gratuit)
        \item Cycles d'achat lents en \'education
        \item D\'ependance aux appareils des \'el\`eves
      \end{itemize}
    \end{block}
  \end{columns}
\end{frame}

%======================================================================
\section{Conclusion}
%======================================================================

\begin{frame}{Conclusion}
  \begin{columns}
    \column{0.55\textwidth}
    \begin{block}{Ce que nous avons r\'ealis\'e}
      \begin{itemize}
        \item[$\checkmark$] Moteur math\'ematique temps r\'eel (\'equation $\to$ courbe)
        \item[$\checkmark$] Physique 2D avec particules, sons, collisions
        \item[$\checkmark$] 8 niveaux p\'edagogiques progressifs
        \item[$\checkmark$] Mode multijoueur LAN via QR Code (Socket.io)
        \item[$\checkmark$] Personnalisation (Garage) avec persistance
        \item[$\checkmark$] Architecture modulaire et maintenable
      \end{itemize}
    \end{block}
    \column{0.42\textwidth}
    \begin{exampleblock}{Technologies ma\^itris\'ees}
      \tech{React 19} $\cdot$ \tech{Canvas API}\\
      \tech{Socket.io} $\cdot$ \tech{Express}\\
      \tech{Web Audio API} $\cdot$ \tech{new Function}\\
      \tech{requestAnimationFrame}\\
      \tech{localStorage} $\cdot$ \tech{QRCode}
    \end{exampleblock}
    \vspace{0.4cm}
    \begin{alertblock}{Prochaine \'etape}
      D\'eploiement sur \textbf{Vercel} (front) + \textbf{Railway} (socket)
      pour des tests en conditions r\'eelles de classe.
    \end{alertblock}
  \end{columns}
\end{frame}

%---------- SLIDE FIN -------------------------------------------------
\begin{frame}[plain]
  \vfill
  \centering
  {\Huge \textcolor{mycyan}{\textbf{GRAPH}}\textbf{RACER}}\\[0.5cm]
  {\Large \textcolor{mylightslate}{Merci pour votre attention !}}\\[1cm]
  {\large \textbf{SAAD OMIT} \quad \& \quad \textbf{IZADINE}}\\[0.5cm]
  \begin{tikzpicture}
    \draw[mycyan, thick, ->] (0,0) -- (7,0);
    \draw[myyellow, very thick, domain=0:6.8, samples=60]
      plot (\x, {0.8*sin(\x r * 0.8) + 0.8});
    \fill[myyellow] (3.0, 1.45) rectangle (3.6, 1.75);
    \fill[darkbg] (3.1, 1.3) circle (0.12);
    \fill[darkbg] (3.5, 1.3) circle (0.12);
    \node[myred] at (6.5, 1.5) {\Large $\blacksquare$};
  \end{tikzpicture}\\[0.4cm]
  {\small \textcolor{mylightslate}{
    \texttt{github.com/saadOmt/GraphRacer}
    \quad $|$ \quad
    Stack: React $\cdot$ Socket.io $\cdot$ Canvas API
  }}
  \vfill
\end{frame}

\end{document}
